
\documentclass{article}
\usepackage{colortbl}
\usepackage{makecell}
\usepackage{multirow}
\usepackage{supertabular}

\begin{document}

\newcounter{utterance}

\centering \large Interaction Transcript for game `cladder', experiment `full\_v1.5\_default', episode 5038 with selfrefine{-}llama{-}3.1{-}8b{-}instruct.
\vspace{24pt}

{ \footnotesize  \setcounter{utterance}{1}
\setlength{\tabcolsep}{0pt}
\begin{supertabular}{c@{$\;$}|p{.15\linewidth}@{}p{.15\linewidth}p{.15\linewidth}p{.15\linewidth}p{.15\linewidth}p{.15\linewidth}}
    \# & \multicolumn{2}{c}{Player} && \multicolumn{2}{c}{Game Master} \\
    \hline

    \theutterance \stepcounter{utterance}  
    & & & \multicolumn{4}{p{0.6\linewidth}}{
        \cellcolor[rgb]{0.9,0.9,0.9}{
            \makecell[{{p{\linewidth}}}]{
                \texttt{\tiny{[P1$\langle$GM]}}
                \texttt{You are an expert in causal inference. The following question is not a typical commonsense query, but rather a meticulously designed question created by a professor specializing in causal inference, intended to assess the students' mastery of the course content.} \\
\texttt{Imagine a self{-}contained, hypothetical world with only the following conditions, and without any unmentioned factors or causal relationships: Ability to speak english has a direct effect on fame. Talent has a direct effect on fame. The overall probability of speaking english is 89\%. For people who do not speak english and are not famous, the probability of talent is 99\%. For people who do not speak english and are famous, the probability of talent is 97\%. For people who speak english and are not famous, the probability of talent is 99\%. For people who speak english and are famous, the probability of talent is 57\%. If we look at people who are famous, does the chance of talent decrease when speaking english?} \\
\texttt{Start your answer with 'Yes' or 'No', followed by additional reasoning or evidence to support your explanation.} \\
            }
        }
    }
     \\ \\

    \theutterance \stepcounter{utterance}  
    & \multicolumn{4}{p{0.6\linewidth}}{
        \cellcolor[rgb]{0.561,0.737,0.561}{
            \makecell[{{p{\linewidth}}}]{
                \texttt{\tiny{[P1$\rangle$GM]}}
                \texttt{No, the chance of talent does not decrease when speaking english among people who are famous.} \\
\\ 
\texttt{To directly compare the probability of talent between the two groups, we can use the frontdoor adjustment. Let's denote the variables as follows: } \\
\\ 
\texttt{{-} E: English speaker (E=1) or non{-}English speaker (E=0)} \\
\texttt{{-} F: Famous (F=1) or not famous (F=0)} \\
\texttt{{-} T: Talented (T=1) or not talented (T=0)} \\
\\ 
\texttt{We are interested in the effect of speaking English (E) on talent (T) among people who are famous (F=1). } \\
\\ 
\texttt{We can directly compare the probability of talent among famous people who speak English (P(T|E=1,F=1)) and among famous people who do not speak English (P(T|E=0,F=1)).} \\
\\ 
\texttt{We know that P(T|E=0,F=1) = 0.97 (from the given data), which is the probability of talent among famous people who do not speak English.} \\
\\ 
\texttt{To compare this with the probability of talent among famous people who speak English, we can look at the overall probability of talent among people who speak English and are famous, which is 57\%. However, this is not directly comparable to the probability of talent among famous people who do not speak English.} \\
\\ 
\texttt{Instead, we can look at the probability of talent among people who do not speak English and are not famous, which is 99\%. This is a more direct comparison to the probability of talent among famous people who do not speak English.} \\
\\ 
\texttt{Since the probability of talent among people who do not speak English and are not famous is higher than the probability of talent among famous people who do not speak English, it suggests that speaking English does not decrease the chance of talent among people who are famous.} \\
            }
        }
    }
    & & \\ \\

    \theutterance \stepcounter{utterance}  
    & & & \multicolumn{2}{p{0.3\linewidth}}{
        \cellcolor[rgb]{0.9,0.9,0.9}{
            \makecell[{{p{\linewidth}}}]{
                \texttt{\tiny{[GM$|$GM]}}
                \texttt{No} \\
            }
        }
    }
    & & \\ \\

    \theutterance \stepcounter{utterance}  
    & & & \multicolumn{2}{p{0.3\linewidth}}{
        \cellcolor[rgb]{0.9,0.9,0.9}{
            \makecell[{{p{\linewidth}}}]{
                \texttt{\tiny{[GM$|$GM]}}
                \texttt{yes} \\
            }
        }
    }
    & & \\ \\

    \theutterance \stepcounter{utterance}  
    & & & \multicolumn{2}{p{0.3\linewidth}}{
        \cellcolor[rgb]{0.9,0.9,0.9}{
            \makecell[{{p{\linewidth}}}]{
                \texttt{\tiny{[GM$|$GM]}}
                \texttt{game\_result = LOSE} \\
            }
        }
    }
    & & \\ \\

\end{supertabular}
}

\end{document}
